% GENERATED FILE. Edit canonical lessons, then run npm run generate:books.
% canonical-source-hash: fnv1a64:0864f857af118280
% canonical-lessons: RU-C189-chuzhoy, RU-C189-rodnoy, RU-C189-mestnyy, RU-C189-inostrannyy, RU-C189-russkiy

\chapter{Someone else's, Native, Local, Foreign, Russian}
\label{ch:ru-someone-else-s-native-local-foreign-russian}
\hlchaptermodality{189}

\begin{chapteropening}
\textbf{By the end of this chapter:} \emph{I can say someone else's, native, local, foreign and Russian.}

Complete the last lesson of chapter 189: I can say someone else's, native, local, foreign and Russian.
\end{chapteropening}

\section[\texorpdfstring{\ru{чужой} (chuzhóy)}{чужой (chuzhóy)}]{\ru{чужой} (chuzhóy) — someone else's; unfamiliar}

\label{lesson:RU-C189-chuzhoy}



\begin{quote}
Before the new one: say the Russian for possible, then the Russian for complicated.
\end{quote}

\subsection*{You'll want to know: \ru{чужой}}
\textbf{\ru{чужой}} (\emph{chuzhóy}) — \textquotedblleft{}someone else's; unfamiliar\textquotedblright{}.

\begin{grammarlens}[title={What you've built}]
One more word for this chapter.
\end{grammarlens}

\subsection*{Your turn}
\begin{itemize}
\raggedright
  \item \emph{Say it:} \emph{chuzhóy}
  \item \emph{Say it:} \emph{chuzhóy}, once more
  \item \emph{Say it:} say \emph{chuzhóy}
  \item \emph{Recall:} say the Russian for possible, then the Russian for complicated, then say \emph{chuzhóy} again
\end{itemize}

\begin{tcolorbox}[breakable,colback=teal!4,colframe=teal!35!black,title={Before you move on}]
What does \ru{чужой} mean? (\textquotedblleft{}Someone else's; unfamiliar\textquotedblright{}.) Say it once more.
\end{tcolorbox}
\section[\texorpdfstring{\ru{родной} (rodnóy)}{родной (rodnóy)}]{\ru{родной} (rodnóy) — native; dear}

\label{lesson:RU-C189-rodnoy}



\begin{quote}
Before the new one: say the Russian for complicated, then the Russian for someone else's; unfamiliar.
\end{quote}

\subsection*{You'll want to know: \ru{родной}}
\textbf{\ru{родной}} (\emph{rodnóy}) — \textquotedblleft{}native; dear\textquotedblright{}.

\begin{grammarlens}[title={What you've built}]
One more word for this chapter.
\end{grammarlens}

\subsection*{Your turn}
\begin{itemize}
\raggedright
  \item \emph{Say it:} \emph{rodnóy}
  \item \emph{Say it:} \emph{rodnóy}, once more
  \item \emph{Say it:} say \emph{rodnóy}
  \item \emph{Recall:} say the Russian for complicated, then the Russian for someone else's; unfamiliar, then say \emph{rodnóy} again
\end{itemize}

\begin{tcolorbox}[breakable,colback=teal!4,colframe=teal!35!black,title={Before you move on}]
What does \ru{родной} mean? (\textquotedblleft{}Native; dear\textquotedblright{}.) Say it once more.
\end{tcolorbox}
\section[\texorpdfstring{\ru{местный} (méstnyy)}{местный (méstnyy)}]{\ru{местный} (méstnyy) — local}

\label{lesson:RU-C189-mestnyy}



\begin{quote}
Before the new one: say the Russian for someone else's; unfamiliar, then the Russian for native; dear.
\end{quote}

\subsection*{You'll want to know: \ru{местный}}
\textbf{\ru{местный}} (\emph{méstnyy}) — \textquotedblleft{}local\textquotedblright{}.

\begin{grammarlens}[title={What you've built}]
One more word for this chapter.
\end{grammarlens}

\subsection*{Your turn}
\begin{itemize}
\raggedright
  \item \emph{Say it:} \emph{méstnyy}
  \item \emph{Say it:} \emph{méstnyy}, once more
  \item \emph{Say it:} say \emph{méstnyy}
  \item \emph{Recall:} say the Russian for someone else's; unfamiliar, then the Russian for native; dear, then say \emph{méstnyy} again
\end{itemize}

\begin{tcolorbox}[breakable,colback=teal!4,colframe=teal!35!black,title={Before you move on}]
What does \ru{местный} mean? (\textquotedblleft{}Local\textquotedblright{}.) Say it once more.
\end{tcolorbox}
\section[\texorpdfstring{\ru{иностранный} (inostránnyy)}{иностранный (inostránnyy)}]{\ru{иностранный} (inostránnyy) — foreign}

\label{lesson:RU-C189-inostrannyy}



\begin{quote}
Before the new one: say the Russian for native; dear, then the Russian for local.
\end{quote}

\subsection*{You'll want to know: \ru{иностранный}}
\textbf{\ru{иностранный}} (\emph{inostránnyy}) — \textquotedblleft{}foreign\textquotedblright{}.

\begin{grammarlens}[title={What you've built}]
One more word for this chapter.
\end{grammarlens}

\subsection*{Your turn}
\begin{itemize}
\raggedright
  \item \emph{Say it:} \emph{inostránnyy}
  \item \emph{Say it:} \emph{inostránnyy}, once more
  \item \emph{Say it:} say \emph{inostránnyy}
  \item \emph{Recall:} say the Russian for native; dear, then the Russian for local, then say \emph{inostránnyy} again
\end{itemize}

\begin{tcolorbox}[breakable,colback=teal!4,colframe=teal!35!black,title={Before you move on}]
What does \ru{иностранный} mean? (\textquotedblleft{}Foreign\textquotedblright{}.) Say it once more.
\end{tcolorbox}
\section[\texorpdfstring{\ru{русский} (rússkiy)}{русский (rússkiy)}]{\ru{русский} (rússkiy) — Russian}

\label{lesson:RU-C189-russkiy}



\begin{quote}
Before the new one: say the Russian for local, then the Russian for foreign.
\end{quote}

\subsection*{You'll want to know: \ru{русский}}
\textbf{\ru{русский}} (\emph{rússkiy}) — \textquotedblleft{}Russian\textquotedblright{}.

\begin{grammarlens}[title={What you've built}]
One more word for this chapter.
\end{grammarlens}

\subsection*{Your turn}
\begin{itemize}
\raggedright
  \item \emph{Say it:} \emph{rússkiy}
  \item \emph{Say it:} \emph{rússkiy}, once more
  \item \emph{Say it:} say \emph{rússkiy}
  \item \emph{Recall:} say the Russian for local, then the Russian for foreign, then say \emph{rússkiy} again
\end{itemize}

\begin{tcolorbox}[breakable,colback=teal!4,colframe=teal!35!black,title={Before you move on}]
What does \ru{русский} mean? (\textquotedblleft{}Russian\textquotedblright{}.) Say it once more.
\end{tcolorbox}
